\clearpage\pagestyle{otMain}\setlength{\otBlocEspacement}{2.5mm}\fontsize{11}{15}\selectfont\corrigesChapitre{16}{Polygones, du triangle à l’hexagone}
\begin{blocLivre}{cor-16-1}\begin{corrige}{1}{Compter les côtés}\demarche{Compter côtés ou sommets.}
Trois côtés : $\res{\text{triangle}}$ ; quatre : $\res{\text{quadrilatère}}$ ; cinq : $\res{\text{pentagone}}$ ; six : $\res{\text{hexagone}}$. On compte les côtés ou les sommets.
\label{fig-cor-16-1}\figureGeo[corrige]{tri-1}
\end{corrige}\end{blocLivre}
\begin{blocLivre}{cor-16-2}\begin{corrige}{2}{Classer les triangles}\demarche{Lire les codes avant l’allure.}
Le petit carré signale le $\res{\text{rectangle}}$ ; deux traits identiques, l’$\res{\text{isocèle}}$ ; trois, l’$\res{\text{équilatéral}}$ ; aucun code : $\res{\text{quelconque}}$.
\label{fig-cor-16-2}\figureGeo[corrige]{tri-2}
\end{corrige}\end{blocLivre}
\begin{blocLivre}{cor-16-3}\begin{corrige}{3}{L’isocèle au compas}\demarche{Même ouverture de compas pour les deux arcs.}
Deux arcs de rayon 6 cm centrés aux extrémités de la base se croisent au sommet principal. Les codes verts affirment les côtés égaux.
\label{fig-cor-16-3}\figureGeo[corrige]{tri-3}
\end{corrige}\end{blocLivre}
\begin{blocLivre}{cor-16-4}\begin{corrige}{4}{Trapèze ou pas ?}\demarche{Contrôler l’écart constant entre côtés opposés.}
Les figures 1 et 3 ont deux côtés opposés $\res{\text{parallèles}}$ : des trapèzes ; la figure 1, avec ses deux angles droits codés, est un trapèze $\res{\text{rectangle}}$. Les figures 2 et 4 n’ont aucune paire de côtés parallèles.
\label{fig-cor-16-4}\figureGeo[corrige]{tri-4}
\end{corrige}\end{blocLivre}
\begin{blocLivre}{cor-16-5}\begin{corrige}{5}{La carte d’identité}\demarche{Lire les codes de la figure.}
Il porte $\res2$ angles droits, posés sur le même côté. Les côtés parallèles sont les deux $\res{\text{bases}}$, la grande et la petite ; les deux autres côtés ne le sont pas.
\label{fig-cor-16-5}\figureGeo[corrige]{tri-5}
\end{corrige}\end{blocLivre}
\begin{blocLivre}{cor-16-6}\begin{corrige}{6}{Le rectangle contrôlé}\demarche{Mesurer les deux diagonales.}
Les deux diagonales mesurent chacune environ $\res{8{,}1}$ cm : elles ont la $\res{\text{même longueur}}$ et se coupent en leur milieu. Ce double contrôle valide la construction.
\label{fig-cor-16-6}\figureGeo[corrige]{tri-6}
\end{corrige}\end{blocLivre}
\begin{blocLivre}{cor-16-7}\begin{corrige}{7}{Le carré incliné}\demarche{S’appuyer sur le côté donné.}
$\res{\text{Non}}$ : perpendiculaires aux extrémités, report de 5 cm, jonction : le carré reste un carré, quelle que soit son inclinaison sur la feuille.
\label{fig-cor-16-7}\figureGeo[corrige]{tri-7}
\end{corrige}\end{blocLivre}
\begin{blocLivre}{cor-16-8}\begin{corrige}{8}{Le losange au compas}\demarche{Même ouverture du compas pour les quatre côtés.}
Depuis le sommet, reporter 4 cm sur chaque côté de l’angle, puis croiser deux arcs de rayon 4 cm pour fermer la figure : quatre côtés égaux, codés en vert.
\label{fig-cor-16-8}\figureGeo[corrige]{tri-8}
\end{corrige}\end{blocLivre}
\begin{blocLivre}{cor-16-9}\begin{corrige}{9}{L’hexagone régulier}\demarche{Ne jamais changer l’ouverture du compas.}
Le rayon reporté six fois sur le cercle donne les six sommets : chaque côté vaut le rayon, 4 cm. Périmètre : $6\times4=\res{24}$ cm. La construction verte garde les arcs.
\label{fig-cor-16-9}\figureGeo[corrige]{tri-9}
\end{corrige}\end{blocLivre}
\begin{blocLivre}{cor-16-10}\begin{corrige}{10}{Régulier ou non ?}\demarche{Vérifier côtés et angles, pas l’allure.}
Le premier a cinq côtés codés égaux et des angles tous pareils : $\res{\text{régulier}}$. Le second a des côtés visiblement inégaux : irrégulier. Le nom « pentagone » dit le nombre de côtés, pas leur égalité.
\label{fig-cor-16-10}\figureGeo[corrige]{tri-10}
\end{corrige}\end{blocLivre}
\begin{blocLivre}{cor-16-11}\begin{corrige}{11}{Le trapèze décomposé}\demarche{Abaisser une perpendiculaire depuis la petite base.}
La perpendiculaire issue du sommet de la petite base découpe un rectangle de $\res{4\times3}$ cm et un triangle rectangle de côtés de l’angle droit $\res{2}$ cm ($6-4$) et $\res3$ cm. Le découpage figure en vert.
\label{fig-cor-16-11}\figureGeo[corrige]{tri-11}
\end{corrige}\end{blocLivre}
\begin{blocLivre}{cor-16-12}\begin{corrige}{12}{Le programme de l’hexagone}\demarche{Utiliser l’égalité côté = rayon.}
Programme : trace un cercle de centre O et de rayon 3 cm ; choisis un point A sur le cercle ; sans changer l’ouverture du compas, reporte le rayon six fois autour du cercle à partir de A ; joins les six points voisins. Le périmètre obtenu vaut $6\times3=\res{18}$ cm.
\label{fig-cor-16-12}\figureGeo[corrige]{tri-12}
\end{corrige}\end{blocLivre}
\begin{blocLivre}{cor-16-13}\begin{corrige}{13}{La table basse}\demarche{Additionner tous les côtés du contour.}
Le tour additionne les quatre côtés : $90+60+40+40=\res{230}$ cm, soit $2{,}30$ m.
\label{fig-cor-16-13}\figureGeo[corrige]{tri-13}
\end{corrige}\end{blocLivre}
\begin{blocLivre}{cor-16-14}\begin{corrige}{14}{Les alvéoles}\demarche{Observer une alvéole entourée de six autres.}
Chaque alvéole entourée partage ses $\res6$ côtés, un avec chaque voisine : aucune paroi n’est construite en double, la cire est $\res{\text{économisée}}$. Le pavage ne laisse aucun trou.
\label{fig-cor-16-14}\figureGeo[corrige]{tri-14}
\end{corrige}\end{blocLivre}
\begin{blocLivre}{cor-16-15}\begin{corrige}{15}{Qui est qui ?}\demarche{Passer chaque figure au crible des définitions.}
« Tout carré est un losange » : $\res{\text{vrai}}$, quatre côtés égaux. « Tout losange est un carré » : $\res{\text{faux}}$, il peut manquer les angles droits. « Tout rectangle a ses côtés opposés parallèles » : $\res{\text{vrai}}$, c’est une propriété du rectangle.
\label{fig-cor-16-15}\figureGeo[corrige]{tri-15}
\end{corrige}\end{blocLivre}
